%% ma: L12-B01-0038
%% lop: 12
%% chuong: 1
%% bai: 1
%% dang: 12.1.24
%% loai: DS
%% muc: [NB, TH, TH, TH]
%% dap_an: "SSĐS"
%% nguon: "HK1-DeOn_1_1.pdf (M10, Đề ôn tập học kì 1 Toán 12, Đề số 2), Phần II Câu 2"
%% kien_thuc: "Bài 1 (tính đơn điệu và cực trị)"
%% trang_thai: cho-duyet
%% ly_do: "Dạng mới đề xuất 12.1.24"
\begin{ex}
Cho hàm số $y=f(x)=x^3+8x^2+5x+1$.
\choiceTF{Đạo hàm $f'(x)=3x^2+8x+5$}{Hàm số đã cho đồng biến trên khoảng $(-1;+\infty)$}{\True Hai điểm cực trị của đồ thị hàm số nằm cùng phía so với trục $Oy$}{$f(0)<f(x)$ với mọi $x\in\mathbb R$}
\loigiai{a) Sai: $f'(x)=3x^2+16x+5$.

b) Sai: $f'(x)=(3x+1)(x+5)<0$ trên $\left(-5;-\dfrac13\right)$, nên hàm số nghịch biến trên $\left(-1;-\dfrac13\right)\subset(-1;+\infty)$.

c) Đúng: hai điểm cực trị có hoành độ $-5$ và $-\dfrac13$, cùng âm.

d) Sai: chẳng hạn $f(-10)=-1000+800-50+1=-249<f(0)=1$ (và tại $x=0$ bất đẳng thức không xảy ra).}
\end{ex}
